\section{Lecture 9: Localization of Solid Analytic Rings}
\label{lec:9}

Clausen turns to \emph{geometry}: how the derived category of an analytic ring
glues over a space of ``points''. The lecture has two halves. First, some
loose ends from Scholze's \cref{lec:8} are tied off --- the existence of the
reflector into $\Mod(R)$, and the stability of the derived category $D(R)$
under arbitrary products --- and the class of \emph{solid} analytic ring
structures $(\ur R,R^{+})^{\solid}$ cut out by a ring of integral elements is
recalled. Second, and the heart of the lecture, is a \emph{descent} theorem:
for a discrete Huber pair $(R,R^{+})$ the assignment
\[
U\ \longmapsto\ D\bigl((\mathcal O(U),\mathcal O^{+}(U))^{\solid}\bigr)
\]
is a sheaf of $\infty$-categories on the \emph{valuative spectrum}
$\Spv(R,R^{+})$, exactly parallel to the Zariski descent of $D(R)$ over
$\Spec R$ --- but only at the derived level, since the localization functors
are the $\mathbb{Z}[T]$-solidifications of \cref{lec:7} and are not $t$-exact.
The proof reduces, by a Tate-style cover-generation lemma, to a single
elementary vanishing of a solid tensor product. The lecture closes by
comparing the whole valuative spectrum with Huber's smaller adic spectrum
$\Spa(R,R^{+})$ of \emph{continuous} valuations, which is recovered as a
retract.

Throughout, ``ring'' means commutative ring, everything lives in light
condensed abelian groups, and we use freely the solid theory
$\Solid\subseteq\CondAbl$ (\cref{lec:5,lec:6}), the relative solid theory over
$\mathbb{Z}[T]$ (\cref{lec:7}), and the abstract notion of analytic ring
(\cref{lec:8}).

\subsection{Recap: analytic rings}

Recall from \cref{lec:8} the notion of an analytic ring: a pair $R=(\ur R,
\Mod(R))$, denoted by the single letter $R$, where $\ur R$ is a condensed ring
(Clausen's triangle notation) and
\begin{equation}
\label{eq:9:analytic-ring}
\Mod(R)\ \subseteq\ \Mod_{\ur R}
\end{equation}
is a full subcategory of the condensed $\ur R$-modules --- the \emph{complete}
modules --- satisfying the closure axioms:
\begin{enumerate}
\item $\Mod(R)$ is closed under all limits, colimits and extensions (in
particular it is an abelian subcategory, but more, since infinite limits and
colimits are allowed);
\item if $M\in\Mod_{\ur R}$ and $N\in\Mod(R)$, then $\iExt^{i}_{\ur R}(M,N)\in
\Mod(R)$ for all $i$ (internal Ext; this is what produces a good tensor product
on $\Mod(R)$ and on its derived analogue);
\item the unit $\ur R$ itself lies in $\Mod(R)$.
\end{enumerate}
The inclusion admits a left adjoint (\cref{prop:9:reflection} below).

The derived category was defined (\cref{def:8:derived-cat}) as the full
subcategory
\begin{equation}
\label{eq:9:derived-cat}
D(R)\ \subseteq\ D(\Mod_{\ur R})\ (=D(\ur R))
\end{equation}
of those $M$ with $H_{i}M\in\Mod(R)$ for all $i$, and one wants $D(R)$ to
inherit the analogues of (1)--(3).

\subsection{Two deferred technical points}

\subsubsection*{Existence of the left adjoint}

The first point left open in \cref{lec:8} was whether the reflector
$\Mod_{\ur R}\to\Mod(R)$ exists automatically. It does, by a version of the
adjoint functor theorem.

\begin{theorem}[Reflection principle; Ad\'amek--Rosick\'y \cite{Adamek_1994}]
\label{prop:9:reflection}
Let $\mathcal C$ be a presentable (\emph{locally presentable}) category and
$\mathcal D\subseteq\mathcal C$ a full subcategory closed under all limits.
Suppose moreover that there exists a regular cardinal $\kappa$ such that
$\mathcal D$ is closed under all $\kappa$-filtered colimits. Then $\mathcal D$
is presentable, and the inclusion $\mathcal D\hookrightarrow\mathcal C$ admits a
left adjoint.
\end{theorem}

Closure under all limits is the condition under which one would naively hope for
a left adjoint; the extra clause on $\kappa$-filtered colimits (for
$\kappa=\aleph_{0}$ this is ordinary filtered colimits, and larger $\kappa$
gives a weaker hypothesis) removes the set-theoretic obstructions. The
$\infty$-categorical version --- needed for $D(R)$ --- is due to
Ragimov--Schlank \cite{ragimov2026inftycategoricalreflectiontheoremapplications}.

\begin{remark}[An alternative for $D(R)$]
\label{rmk:9:presentable-derived}
For the derived category one can avoid invoking the $\infty$-categorical
reflection theorem. There is a general principle (attributed in discussion to
Scholze) that presentable categories are closed under limits inside the
category of (large) categories, provided the connecting functors preserve
colimits; since the homology functors $H_{i}$ preserve colimits, $D(R)$ --- cut
out by conditions on the $H_{i}$ --- is presentable, and then Lurie's more
naive adjoint functor theorem \cite{lurie2017higher} supplies the left adjoint.
One can also check presentability at the level of complexes directly. Clausen
was content to believe the $\infty$-categorical input is avoidable.
\end{remark}

\subsubsection*{Stability of $D(R)$ under products}

The second point was closure of $D(R)$ under limits, which reduces to closure
under products. The subtlety (\cref{lec:3}) is that in light condensed abelian
groups \emph{countable} products are exact but arbitrary products are not, so
the product functor has higher right derived functors $R^{i}\!\prod$, and one
must check these land in $\Mod(R)$.

\begin{proposition}[Derived products stay complete]
\label{prop:9:products}
Let $(M_{\alpha})_{\alpha\in A}$ be a family of objects of $\Mod(R)$. Viewing
each $M_{\alpha}$ in $D(R)$ in degree $0$, the derived product
$\prod_{\alpha}M_{\alpha}$ lies in $D(R)$; equivalently
$R^{i}\!\prod_{\alpha\in A}M_{\alpha}\in\Mod(R)$ for all $i$.
\end{proposition}

\begin{proof}
Set $M=\bigoplus_{\alpha\in A}M_{\alpha}\in\Mod(R)$. Termwise, the system
$(M_{\alpha})$ is a retract of the constant-in-$M$ system, so
$R^{i}\!\prod_{\alpha}M_{\alpha}$ is a retract of $R^{i}\!\prod_{\alpha\in A}M$.
But the latter is internal Ext out of a free module,
\begin{equation}
\label{eq:9:product-as-ext}
R^{i}\!\prod_{\alpha\in A}M\ =\ \iExt^{i}_{\ur R}\Bigl(\,\bigoplus_{\alpha\in A}
\ur R,\ M\,\Bigr),
\end{equation}
which lies in $\Mod(R)$ by axiom~(2). Since $\Mod(R)$ is closed under retracts
(being closed under kernels and cokernels), the claim follows.
\end{proof}

The case of unbounded-below complexes was already handled in \cref{lec:8}:
products there are computed through the countable Postnikov limits, reducing to
the countable-product case.

\begin{remark}[Light vs.\ general condensed]
\label{rmk:9:light-vs-general}
Light condensed rings embed fully faithfully into all condensed rings, so the
first coordinate transports without loss. The analytic ring \emph{structure} is
more subtle: a light analytic ring structure only decides what the free modules
should be on \emph{light} profinite sets, so passing to the general condensed
setting is not a mere restriction of scalars.
\end{remark}

\subsection{Solid analytic ring structures}

We specialise to the solid context. A \emph{solid ring} $\ur R$ is a condensed
ring whose underlying condensed abelian group is solid, i.e.\ a commutative
algebra object
\begin{equation}
\label{eq:9:solid-ring}
\ur R\in\CAlg(\Solid_{\mathbb{Z}}),
\end{equation}
and it carries the subset $R^{\circ}\subseteq\ur R(\ast)$ of power-bounded
elements (\cref{def:7:tn-pb}).

\begin{definition}[The structure $(\ur R,R^{+})^{\solid}$]
\label{def:9:solid-structure}
For any subset $R^{+}\subseteq R^{\circ}$ with $1\in R^{+}$, one obtains an
analytic ring structure on $\ur R$, denoted $(\ur R,R^{+})^{\solid}$, with the
same underlying ring and complete modules
\begin{equation}
\label{eq:9:solid-mod}
\Mod\bigl((\ur R,R^{+})^{\solid}\bigr)\ =\ \Bigl\{\, M\in\Mod_{\ur R}\ \Bigm|\
\iHom(P,M)\xrightarrow{\ 1-\sigma f\ }\iHom(P,M)\ \text{is an iso}\ \forall f\in
R^{+}\,\Bigr\},
\end{equation}
where $P=\mathbb{Z}[\mathbb{N}\cup\{\infty\}]/\mathbb{Z}[\infty]$ is the free
group on a null sequence and $\sigma$ its shift (\cref{lec:5,lec:7}). In words,
one requires that for each $f\in R^{+}$ the module be $\mathbb{Z}[T]$-solid
along $T\mapsto f$.
\end{definition}

\begin{remark}[Reading the two conditions]
\label{rmk:9:reading-conditions}
The two boundary cases $f=1$ and $R^{+}\subseteq R^{\circ}$ have clean meanings.
\begin{itemize}
\item Taking $f=1$ recovers the definition of $\Solid_{\mathbb{Z}}$
(\cref{def:5:solid}); so $1\in R^{+}$ forces every $M\in\Mod(R)$ to be solid as
a condensed abelian group.
\item The condition $R^{+}\subseteq R^{\circ}$ (power-bounded) is precisely what
makes the ring $\ur R$ itself complete, i.e.\ $\ur R\in\Mod(R)$ (axiom~(3)):
the second coordinate specifies a completeness notion adapted to the situation,
and one wants the ring to be already complete.
\end{itemize}
\end{remark}

\begin{remark}[No loss in taking $R^{+}$ integrally closed]
\label{rmk:9:integrally-closed}
There is no loss of generality in assuming $R^{+}\subseteq\ur R(\ast)$ is an
integrally closed subring containing all topologically nilpotent elements
$R^{\circ\circ}$. Indeed, given any solid analytic ring structure on $\ur R$
(one for which every module is solid), the set of $f$ for which
$1-\sigma f$ acts invertibly on $\iHom(P,M)$ for all such $M$ is already an
integrally closed subring containing $R^{\circ\circ}$
(\cref{lem:7:Rzero,prop:8:Aplus-integrally-closed}). So one may always replace
$R^{+}$ by the integral closure of the subring it generates without changing the
theory.
\end{remark}

\begin{example}[Huber pairs]
\label{ex:9:huber-pair}
If $\ur R$ is (the condensed ring associated to) a complete Huber ring, then the
integrally closed subrings $R^{+}$ with $R^{\circ\circ}\subseteq R^{+}\subseteq
R^{\circ}$ are the same as the \emph{open} integrally closed subrings
$R^{+}\subseteq R^{\circ}$ --- exactly the rings of integral elements of Huber's
theory. So the general classification of admissible $R^{+}$ specialises,
for a Huber ring, to Huber's own. Moreover $R^{+}$ is \emph{recovered} from the
analytic ring $(\ur R,R^{+})^{\solid}$: the assignment
\begin{equation}
\label{eq:9:huber-embedding}
(\ur R,R^{+})\ \longmapsto\ (\ur R,R^{+})^{\solid}
\end{equation}
is fully faithful, so a Huber pair and its solid analytic ring carry the same
information (\cref{prop:8:huber-fully-faithful,thm:8:correspondence};
the converse --- that no \emph{other} $f$ can satisfy the defining condition ---
was proved for Huber rings in \cref{lec:8}). ``Huber ring'' henceforth means
\emph{complete} Huber ring.
\end{example}

\begin{remark}[Derived-complete rings]
\label{rmk:9:derived-complete}
Asked about ``derived complete'' rings, Clausen confirmed these too give
condensed/analytic objects and can be regarded as solid; there is ``a fun
story'' there, deferred.
\end{remark}

\subsection{Localization: reduction to the discrete case}

\begin{remark}[All the data is already discrete]
\label{rmk:9:discrete-data}
Here is a remark that is shocking at first glance but trivial. Fix a solid ring
$\ur R$ and $R^{+}\subseteq R^{\circ}$ as above (say integrally closed,
containing $R^{\circ\circ}$). The condition \eqref{eq:9:solid-mod} defining the
analytic ring structure is imposed only for $f\in R^{+}$, and $R^{+}$ is a
subset of the \emph{underlying discrete ring} $\ur R(\ast)$. So all the data
used to define the structure already appears at the discrete level: the same
$R^{+}$ defines a structure on the \emph{discrete} ring $\ur R(\ast)$, where
every element is power-bounded, giving the discrete Huber pair
$(\ur R(\ast),R^{+})$.

Consequently, to understand $\Mod\bigl((\ur R,R^{+})^{\solid}\bigr)$ it suffices
to understand the theory over the discrete pair and then remember that $\ur R$
is a commutative algebra object there:
\begin{equation}
\label{eq:9:base-change-from-discrete}
\Mod\bigl((\ur R,R^{+})^{\solid}\bigr)\ =\ \Mod_{\ur R}\Bigl(\Mod\bigl(
(\ur R(\ast),R^{+})^{\solid}\bigr)\Bigr),
\end{equation}
i.e.\ the general theory is base-changed from the discrete case in a completely
naive way. (Because we are doing \emph{condensed} modules, a discrete ring
already has a huge supply of non-discrete modules; in particular $\ur R$ itself
is an honest condensed ring over the discrete condensed ring $\ur R(\ast)$.)
Thus, to understand how these categories glue it is enough to treat the discrete
case, and then put the $\ur R$-module structure on top.
\end{remark}

We describe one framework for gluing --- not the most general, but quite
general. It is organised by two analogies, one with the Zariski spectrum
(classical) and one with the valuative spectrum (the solid analogue).

\subsubsection*{The classical picture: $D(R)$ over $\Spec R$}

\begin{remark}[Zariski localization of $D(R)$]
\label{rmk:9:spec-analogy}
Let $R$ be an ordinary commutative ring and $D(R)$ its usual (unbounded)
derived category of discrete $R$-modules. Then $D(R)$ \emph{localizes over}
$\Spec R=\{\text{prime ideals }\mathfrak p\subseteq R\}$. This space has a basis
of quasicompact opens, closed under finite intersection, the
\emph{distinguished opens}
\begin{equation}
\label{eq:9:distinguished-open}
U(f)\ =\ \{\mathfrak p:\ f\notin\mathfrak p\}\qquad(f\in R),
\end{equation}
carrying a structure sheaf with $\mathcal O(U(f))=R[\tfrac1f]$. Neither $U(f)$
nor $R[\tfrac1f]$ determines $f$, but they determine each other:
$U(f)\cong\Spec R[\tfrac1f]$, matching up distinguished opens. Assigning to a
distinguished open $U$ the derived category $D(\mathcal O(U))$, with base-change
functors for inclusions $V\subseteq U$, gives a presheaf of $\infty$-categories.
\end{remark}

\begin{theorem}[Lurie \cite{lurie2017higher}]
\label{thm:9:zariski-sheaf}
The presheaf $U\mapsto D(\mathcal O(U))$ on distinguished opens of $\Spec R$ is a
sheaf of $\infty$-categories. It is even a sheaf (indeed a hypersheaf) for the
hypercover topology, and there is a sheaf on the level of abelian categories
$U\mapsto\Mod_{\mathcal O(U)}$ as well.
\end{theorem}

\begin{remark}[The abelian sheaf fails in the solid setting]
\label{rmk:9:no-abelian-sheaf}
In the Zariski case one has a sheaf already on abelian categories because the
localization maps $\mathcal O(U)\to\mathcal O(V)$ are \emph{flat}, so derived
statements descend to abelian ones. In the solid analogue below, the
localization maps are \emph{not} flat --- one of them is the
$\mathbb{Z}[T]$-solidification of \cref{lec:7}, which is off from $t$-exactness
by one (\cref{rmk:7:not-t-exact}) --- and this obstructs the sheaf property at
the abelian level. \emph{We obtain a sheaf only on derived categories.}
\end{remark}

\subsubsection*{The solid picture: the valuative spectrum}

Now let $(R,R^{+})$ be a discrete Huber pair: $R$ an ordinary commutative ring
and $R^{+}\subseteq R$ an integrally closed subring. To it we attached the
analytic ring $(R,R^{+})^{\solid}$ and the derived category $D\bigl((R,R^{+})^{
\solid}\bigr)$, abbreviated $D(R,R^{+})$. The claim is that $D(R,R^{+})$
localizes over the \emph{valuative spectrum} $\Spv(R,R^{+})$.

\begin{definition}[Valuative spectrum]
\label{def:9:valuative-spectrum}
Whereas a prime ideal records only whether a function vanishes, a valuation
records the finer comparison of sizes of functions. Define
\begin{equation}
\label{eq:9:spv}
\Spv(R,R^{+})\ =\ \Bigl\{\, v\colon R\longrightarrow\Gamma\cup\{0\}\ \Bigm|\
(\text{axioms below})\,\Bigr\}\big/\!\sim,
\end{equation}
where $\Gamma$ is an ordered abelian group written multiplicatively, and $v$
satisfies:
\begin{align}
\label{eq:9:valuation-axioms}
v(fg)&=v(f)v(g), & v(f+g)&\le\max\bigl(v(f),v(g)\bigr), \notag\\
v(0)&=0,\quad v(1)=1, & v(f)&\le 1\quad\forall f\in R^{+}.
\end{align}
Two valuations are equivalent, $v\sim w$, iff they induce the same order
relation on functions:
\begin{equation}
\label{eq:9:valuation-equivalence}
v\sim w\ \iff\ \bigl(\ \forall f,g\in R:\ v(f)\le v(g)\ \Longleftrightarrow\
w(f)\le w(g)\ \bigr).
\end{equation}
The genuinely important datum of a valuation is thus the binary relation
``$f$ is smaller than $g$''.
\end{definition}

\begin{remark}[Second description; higher rank]
\label{rmk:9:spv-second}
Equivalently, a point of $\Spv(R,R^{+})$ is a pair of a prime ideal
$\mathfrak p\subseteq R$ and a valuation subring (valuation domain) of the
residue field $\kappa(\mathfrak p)$ containing the image of $R^{+}$: the prime
records vanishing, the valuation ring records which fractions of non-vanishing
functions are $\le 1$. Since the continuity condition is vacuous for a discrete
ring, this is the same as Huber's $\Spa$ in the discrete case. For $\mathbb Q$
(or $\mathbb Z$) one has the trivial valuations and the $p$-adic valuations, and
the classification is manageable; but adding one variable makes valuations
proliferate --- in particular \emph{higher-rank} valuations appear (many curves
through a point of a surface), a genuine additional complication.
\end{remark}

\begin{definition}[Rational opens and their structure sheaves]
\label{def:9:rational-open}
The space $\Spv(R,R^{+})$ has a basis of quasicompact opens, closed under finite
intersection, the \emph{rational opens}. For $f_{1},\dots,f_{n},g\in R$,
\begin{equation}
\label{eq:9:rational-open}
U\Bigl(\tfrac{f_{1},\dots,f_{n}}{g}\Bigr)\ =\ \bigl\{\,v:\ v(g)\neq 0,\ v(f_{i})\le
v(g)\ \forall i\,\bigr\}\ \subseteq\ \{v:\ v(g)\neq 0\},
\end{equation}
i.e.\ one sits inside the Zariski open where $g$ is invertible and shrinks by the
inequalities $|f_{i}|\le|g|$. It carries \emph{two} structure sheaves: the
algebraic (Zariski) localization
\begin{equation}
\label{eq:9:structure-sheaf}
\mathcal O\Bigl(U\bigl(\tfrac{f_{1},\dots,f_{n}}{g}\bigr)\Bigr)\ =\ R[\tfrac1g],
\end{equation}
together with a ring of integral elements, the integral closure
\begin{equation}
\label{eq:9:structure-sheaf-plus}
\mathcal O^{+}\Bigl(U\bigl(\tfrac{f_{1},\dots,f_{n}}{g}\bigr)\Bigr)\ =\
\overline{R^{+}\bigl[\tfrac{f_{1}}{g},\dots,\tfrac{f_{n}}{g}\bigr]}\ \subseteq\
R[\tfrac1g],
\end{equation}
obtained by adjoining the $f_{i}/g$ (forced $\le 1$ on the open) to $R^{+}$ and
taking the integral closure. Then
\begin{equation}
\label{eq:9:rational-recursive}
U\Bigl(\tfrac{f_{1},\dots,f_{n}}{g}\Bigr)\ \cong\ \Spv\bigl(\mathcal O(U),
\mathcal O^{+}(U)\bigr),
\end{equation}
matching up rational opens, exactly as $U(f)\cong\Spec R[\tfrac1f]$ in the
Zariski case.
\end{definition}

\begin{remark}[Why $R^{+}$ is needed]
\label{rmk:9:why-Rplus}
Equation \eqref{eq:9:rational-recursive} is where the necessity of carrying the
datum $R^{+}$ becomes visible. Had one dropped the integral condition to get a
bigger space, the rational opens would no longer be of the form
$\Spv(\text{ring pair})$: the abstract ring $R[\tfrac1g]$ does not remember that
the $f_{i}/g$ were forced $\le 1$. To keep rational opens of the same shape as
the global object, one must include $R^{+}$ from the start.
\end{remark}

We assign to a rational open $U$ the derived category $D\bigl((\mathcal O(U),
\mathcal O^{+}(U))^{\solid}\bigr)$, and to an inclusion $U\supseteq V$ the
pullback (base-change) functor
\begin{equation}
\label{eq:9:pullback-functor}
D\bigl((\mathcal O(U),\mathcal O^{+}(U))^{\solid}\bigr)\ \longrightarrow\
D\bigl((\mathcal O(V),\mathcal O^{+}(V))^{\solid}\bigr),
\end{equation}
which comes from a map of analytic rings $(\mathcal O(U),\mathcal O^{+}(U))^{
\solid}\to(\mathcal O(V),\mathcal O^{+}(V))^{\solid}$ (\cref{def:8:maps-analytic-rings}):
restriction of scalars preserves completeness, and its left adjoint --- tensor
up the underlying module, then recomplete --- is the base change.

\begin{theorem}[Descent for the solid affinoid theory]
\label{thm:9:solid-sheaf}
For a discrete Huber pair $(R,R^{+})$, the presheaf on rational opens of
$\Spv(R,R^{+})$,
\begin{equation}
\label{eq:9:solid-presheaf}
U\ \longmapsto\ D\bigl((\mathcal O(U),\mathcal O^{+}(U))^{\solid}\bigr),
\end{equation}
is a sheaf of $\infty$-categories. \emph{Warning:} unlike the classical case, it
is \emph{not} a sheaf at the level of abelian categories, because the pullback
functors \eqref{eq:9:pullback-functor} are not $t$-exact in general
(\cref{rmk:9:no-abelian-sheaf}).
\end{theorem}

Since the rational opens form a basis closed under finite intersection, a sheaf
on them extends uniquely to a sheaf on all opens (an arbitrary open is covered
by rationals and one takes the limit); so \cref{thm:9:solid-sheaf} determines a
category for every open, rational or not.

\begin{remark}[The discrepancy is bounded]
\label{rmk:9:bounded-discrepancy}
The failure of $t$-exactness is controlled: the $\mathbb{Z}[T]$-solidification is
given by $R\iHom$ against an object with a two-term compact resolution, so it is
off from $t$-exact by at most one --- the homology of the solidification of an
object concentrated in degree $0$ vanishes above degree $n$ for some finite $n$
(indeed $n=1$). The sheaf need not be hypercomplete in general, though
hypercompleteness is automatic when the base space has finite Krull dimension.
\end{remark}

\subsection{Proof of the classical theorem}

To motivate the proof of \cref{thm:9:solid-sheaf} we give a particular proof of
Zariski descent (\cref{thm:9:zariski-sheaf}) that transports with little change.
Work on the site of distinguished opens $U\subseteq\Spec R$ with the open-cover
topology.

\begin{lemma}[Cover-generation]
\label{lem:9:zariski-covers}
The Grothendieck topology on distinguished opens is generated by:
\begin{itemize}
\item the empty cover of the empty set; and
\item for each distinguished open $U\subseteq\Spec R$ and each $f\in\mathcal
O(U)$, the two-element cover
\begin{equation}
\label{eq:9:zariski-basic-cover}
\bigl\{\,U(f),\ U(1-f)\,\bigr\}\ \longrightarrow\ U.
\end{equation}
\end{itemize}
\end{lemma}

\begin{proof}
A general cover of $U$ is given by $f_{1},\dots,f_{n}\in\mathcal O(U)$
generating the unit ideal --- $\exists\,x_{1},\dots,x_{n}$ with $\sum_{i}x_{i}
f_{i}=1$ --- and equals $\{U(f_{i})\}_{i}$. (The empty cover of the empty set
must be listed separately: it cannot be generated from nonempty data, and
checking the sheaf condition there just says the value on $\emptyset$ is
terminal.) Refining $f_{i}\rightsquigarrow f_{i}x_{i}$ we may assume
$\sum_{i}f_{i}=1$, and then induct on $n$. For instance with $f_{1}+f_{2}+f_{3}=
1$, the basic cover $\{U(f_{1}+f_{2}),U(f_{3})\}$ (of the shape
\eqref{eq:9:zariski-basic-cover}, since $f_{3}=1-(f_{1}+f_{2})$) reduces the
sheaf condition to its restrictions: over $U(f_{1}+f_{2})$ the element
$f_{1}+f_{2}$ becomes a unit and one divides through to reach the two-variable
case; over $U(f_{3})$ the cover splits (one of the covering opens is all of it),
so the sheaf condition is automatic. This collection is closed under base
change --- an important technical point --- and that makes the induction go
through. This reduction is originally due to Quillen (in his proof of the
Serre conjecture on projective modules).
\end{proof}

So it suffices to check the sheaf condition for the cover $\{U(f),U(1-f)\}$ of
$U=\Spec R$ (no loss in taking $U=\Spec R$). Concretely this asks that the
functor
% board 22 @ 00:59:18--01:05:58
\begin{equation}
\label{eq:9:zariski-pullback}
D(R)\ \xrightarrow{\ \simeq\ }\ D\bigl(R[\tfrac1f]\bigr)\ \times_{\,D(R[\frac1{
f(1-f)}])}\ D\bigl(R[\tfrac1{1-f}]\bigr)
\end{equation}
be an equivalence of $\infty$-categories, where the right-hand side is the
pullback of the cospan
\begin{equation}
\label{eq:9:zariski-cospan}
\begin{tikzcd}[column sep=large]
& D\bigl(R[\tfrac1f]\bigr) \arrow[d] \\
D\bigl(R[\tfrac1{1-f}]\bigr) \arrow[r] & D\bigl(R[\tfrac1{f(1-f)}]\bigr).
\end{tikzcd}
\end{equation}

\begin{remark}[What the pullback category is, concretely]
\label{rmk:9:pullback-concrete}
An object of the pullback is a triple $(M,N,\alpha)$ of $M\in D(R[\tfrac1f])$,
$N\in D(R[\tfrac1{1-f}])$, and an isomorphism $\alpha\colon M[\tfrac1{1-f}]
\xrightarrow{\sim}N[\tfrac1f]$ of their common base changes --- but an
isomorphism in the \emph{$\infty$-categorical} sense, e.g.\ (for bounded-above
complexes of projectives) an explicit chain-homotopy equivalence, not merely an
iso in the ordinary derived category. Essential surjectivity of
\eqref{eq:9:zariski-pullback} is the statement that such glueing data descends
uniquely to a global object; full faithfulness is that global $R\Homop$'s can be
computed locally, as the homotopy pullback of the three local $R\Homop$-complexes
(a shifted mapping cone). So \eqref{eq:9:zariski-pullback} both glues locally
defined objects and computes global Ext by localizing.
\end{remark}

\begin{proof}[Proof of \eqref{eq:9:zariski-pullback}]
Each base change $D(R)\to D(R[\tfrac1f])$, etc., has a right adjoint (the
forgetful functor), so the comparison functor
\eqref{eq:9:zariski-pullback} has a right adjoint too --- an automatic candidate
inverse. Explicitly it sends $(M,N,\alpha)$ to the fibre product $M\times_{M[
\frac1f]=N[\frac1{1-f}]}N$ (apply the right adjoints and take the limit). One
must check the unit and counit are isomorphisms; both are maps in one of the two
categories. For the unit one needs, for $M\in D(R)$, that the square
% board 24 @ 01:08:25--01:11:10
\begin{equation}
\label{eq:9:unit-square}
\begin{tikzcd}[column sep=large, row sep=large]
M \arrow[r] \arrow[d] & M[\tfrac1f] \arrow[d] \\
M[\tfrac1{1-f}] \arrow[r] & M[\tfrac1{f(1-f)}]
\end{tikzcd}
\end{equation}
is a (homotopy) pullback. The difference between $M$ and the pullback is a
mapping cone $N$ with $N[\tfrac1f]=N[\tfrac1{1-f}]=0$, and the point is that
\begin{equation}
\label{eq:9:conservativity}
N\in D(R),\ \ N[\tfrac1f]=N[\tfrac1{1-f}]=0\ \Longrightarrow\ N=0.
\end{equation}
This reduces to the abelian level (an object vanishes iff its homology does) and,
the localizations being flat, is elementary there. The counit reduces to the
same statement \eqref{eq:9:conservativity}.
\end{proof}

\begin{remark}[The three formal ingredients]
\label{rmk:9:three-ingredients}
The proof uses only three facts about the localizations $M\mapsto M[\tfrac1f]$,
$M\mapsto M[\tfrac1{1-f}]$:
\begin{enumerate}
\item each base change is a \emph{localization} (its right adjoint is fully
faithful);
\item the localizations \emph{commute}: something on which $1-f$ acts invertibly
still has $1-f$ invertible after inverting $f$, and the composite localization
$M\mapsto M[\tfrac1{f(1-f)}]$ is the common further localization;
\item they are \emph{jointly conservative}: if $M\mapsto 0$ under each element of
the cover, then $M=0$.
\end{enumerate}
Given a square of localizations with these three properties, one automatically
gets a pullback. This is what will be verified in the solid setting.
\end{remark}

\subsection{Proof of the solid analogue}

Work on the site of rational opens $U\subseteq\Spv(R,R^{+})$ with the open-cover
topology. The cover-generation lemma now has two families, reflecting the extra
inequalities.

\begin{lemma}[Cover-generation, cf.\ Tate \cite{Tate_1971}]
\label{lem:9:solid-covers}
The topology on rational opens of $\Spv(R,R^{+})$ is generated by:
\begin{itemize}
\item the empty cover of the empty set; and
\item for each rational open $U$ and each $f\in\mathcal O(U)$, the two covers
\begin{equation}
\label{eq:9:solid-covers}
\bigl\{\,U\bigl(\tfrac{f}{1}\bigr),\ U\bigl(\tfrac1f\bigr)\,\bigr\}\quad\text{and}
\quad\bigl\{\,U\bigl(\tfrac1f\bigr),\ U\bigl(\tfrac1{1-f}\bigr)\,\bigr\}.
\end{equation}
\end{itemize}
Here $U(\tfrac f1)=\{v(f)\le 1\}$ is where $f$ is integral (the closed unit
disk); $U(\tfrac1f)=\{v(f)\neq 0,\ v(f)\ge 1\}$ is where $f$ is well away from
$0$ (the complement of the open unit disk). Both are refinements of the Zariski
cover $\{U(f),U(1-f)\}$ that still cover.
\end{lemma}

For the general reduction --- Huber shows every cover is refined by one of the
form: take $f_{1},\dots,f_{n}$ generating the unit ideal, and take the rational
opens $U\bigl(\tfrac{f_{1},\dots,\widehat{f_{i}},\dots,f_{n}}{f_{i}}\bigr)$
(a refinement of the Zariski cover $\{U(f_{i})\}$ that still covers the valuative
spectrum). Then, exactly as in the classical proof, one plays the same reduction
until the theorem reduces to a single case.

\begin{proposition}[The two pullback squares]
\label{prop:9:solid-squares}
For $(R,R^{+})$ a discrete Huber pair and $f\in R$, the theorem reduces to the
statements that
\begin{equation}
\label{eq:9:solid-pullback-1}
D(R,R^{+})\ \xrightarrow{\ \simeq\ }\ D\bigl(R,\ \widetilde{R^{+}[f]}\bigr)\
\times_{\,D(R[\frac1f],\,\widetilde{R^{+}[\frac1f,f]})}\ D\bigl(R[\tfrac1f],\,
\widetilde{R^{+}[\tfrac1f]}\bigr)
\end{equation}
and the analogous square for the second cover in \eqref{eq:9:solid-covers} are
pullbacks, where $\widetilde{(-)}$ denotes integral closure.
\end{proposition}

The identification of the base-change functors is the crux, and here the theory
of \cref{lec:7} enters.

\begin{lemma}[Base changes as $R\iHom$'s]
\label{lem:9:base-change-rhom}
Let $f\in R$.
\begin{enumerate}
\item The closed-disk base change
\begin{equation}
\label{eq:9:closed-disk-map}
D(R,R^{+})\longrightarrow D\bigl(R,\widetilde{R^{+}[f]}\bigr)
\end{equation}
is simply the $\mathbb{Z}[T]$-solidification for $\mathbb{Z}[T]\xrightarrow{\,
T\mapsto f\,}R$, given (\cref{prop:7:solidification-formula}) by
% board 32--33 @ 01:22:40--01:28:53
\begin{equation}
\label{eq:9:closed-disk-rhom}
R\iHom_{\mathbb{Z}[T]}\Bigl(\bigl(\mathbb{Z}(\!(T^{-1})\!)/\mathbb{Z}[T]\bigr)
[-1],\ -\Bigr).
\end{equation}
\item The complementary base change
\begin{equation}
\label{eq:9:complement-map}
D(R,R^{+})\longrightarrow D\bigl(R[\tfrac1f],\widetilde{R^{+}[\tfrac1f]}\bigr)
\end{equation}
--- first invert $f$, then solidify with respect to $1/f$ --- is \emph{also} an
$R\iHom$: for $\mathbb{Z}[T]\xrightarrow{\,T\mapsto f\,}R$,
\begin{equation}
\label{eq:9:complement-rhom}
R\iHom_{\mathbb{Z}[T]}\Bigl(\bigl(\mathbb{Z}[\![T]\!]/\mathbb{Z}[T]\bigr)[-1],\
-\Bigr),
\end{equation}
the localization killing the idempotent object $\mathbb{Z}[\![T]\!]\in
D\bigl(\Mod_{\mathbb{Z}[T]}(\Solid_{\mathbb{Z}})\bigr)$.
\end{enumerate}
\end{lemma}

\begin{proof}[Why \eqref{eq:9:complement-rhom} inverts $f$]
By \cref{lem:7:idempotent} the object $\mathbb{Z}[\![T]\!]$ is idempotent over
$\mathbb{Z}[T]$, so $R\iHom$ against $(\mathbb{Z}[\![T]\!]/\mathbb{Z}[T])[-1]$ is
the Bousfield localization killing $\mathbb{Z}[\![T]\!]$. Killing
$\mathbb{Z}[\![T]\!]$ kills every module over it, and everything $T$-torsion is a
$\mathbb{Z}[\![T]\!]$-module: a solid abelian group that is a filtered colimit of
things killed by powers of $T$ is a $\mathbb{Z}[\![T]\!]$-module (reduce, using
closure under limits and colimits, to something uniformly killed by a power of
$T$, hence a module over the truncated power-series ring). So this localization
automatically inverts $f$. After the change of variables inverting $T$, what
remains is exactly the $1/f$-solidification, i.e.\ the composite ``invert $f$,
then solidify''.
\end{proof}

\begin{proof}[Proof of \cref{thm:9:solid-sheaf}]
By \cref{rmk:9:three-ingredients} it remains to verify the three ingredients for
the two covers. Ingredients (1) and (2) are immediate: each base change is a
localization by construction, and any two localizations given by $R\iHom$ out of
fixed objects commute (as $R\iHom(A,R\iHom(B,-))=R\iHom(A\otimes B,-)$ and the
tensor product is symmetric). Only the joint conservativity (3) must be checked,
and it translates into a single tensor-product vanishing per cover.

For the first cover $\{U(\tfrac f1),U(\tfrac1f)\}$ --- localize to the closed
unit disk, resp.\ away from the open unit disk at $0$ --- ingredient (3) is:
% board 34, 40 @ 01:30:21--01:32:28
\begin{equation}
\label{eq:9:tensor-cover-1}
\mathbb{Z}[\![T]\!]\ \otimes^{\mathbb{L}\solid}_{\mathbb{Z}[T]}\
\mathbb{Z}(\!(T^{-1})\!)\ =\ \mathbb{Z}[\![T,U]\!]/(1-UT)\ =\ 0
\qquad(U=T^{-1}).
\end{equation}
For the second cover $\{U(\tfrac1f),U(\tfrac1{1-f})\}$ --- away from the open
unit disks at $0$ and at $1$ --- it is:
\begin{equation}
\label{eq:9:tensor-cover-2}
\mathbb{Z}[\![T]\!]\ \otimes^{\mathbb{L}\solid}_{\mathbb{Z}[T]}\
\mathbb{Z}[\![1-T]\!]\ =\ \mathbb{Z}[\![T,U]\!]/\bigl(1-(T+U)\bigr)\ =\ 0
\qquad(U=1-T).
\end{equation}
In each case both variables are topologically nilpotent, so the displayed
denominator is $1$ minus a topologically nilpotent element, hence a unit (invert
by geometric series), and the quotient is $0$.
\end{proof}

\begin{remark}[Geometric meaning of the vanishings]
\label{rmk:9:disjoint-disks}
The two identities say that the two idempotent algebras being killed have
disjoint support. Equation \eqref{eq:9:tensor-cover-1} says the open unit disk at
$0$ and the open unit disk at $\infty$ do not meet: their union of complements
(closed disk at $0$ together with closed disk at $\infty$) is the whole space.
Equation \eqref{eq:9:tensor-cover-2} says the open unit disk at $0$ and the open
unit disk at $1$ do not meet --- strange until one recalls this is
non-archimedean geometry, where it is exactly right.
\end{remark}

\subsection{Consequences and comparison with Huber}

\begin{corollary}[Localization for a general solid ring]
\label{cor:9:general-solid}
Let $\ur R$ be any solid ring and $R^{+}\subseteq R^{\circ}\subseteq\ur R(\ast)$
with $1\in R^{+}$. Then $D(\ur R,R^{+})$ localizes on the valuative spectrum
$\Spv(\ur R(\ast),R^{+})$ of the underlying discrete ring: to a rational open
$U$ one assigns
\begin{equation}
\label{eq:9:general-localization}
U\ \longmapsto\ \Mod_{\ur R}\Bigl(D\bigl((\mathcal O(U),\mathcal O^{+}(U))
\bigr)\Bigr),
\end{equation}
the $\ur R$-modules in the (discrete) affinoid category of the rational open.
The unit object over $U(\tfrac{f_{1},\dots,f_{n}}{g})$ is $\ur R[\tfrac1g]$
derived-solidified with respect to the $f_{i}/g$.
\end{corollary}

This is the base-change of the discrete descent (\cref{thm:9:solid-sheaf}) along
$\ur R$, by \cref{rmk:9:discrete-data}. Note that for a nondiscrete $\ur R$ this
unit object generically carries genuine \emph{derived} phenomena (higher
homotopy): the structure sheaf is a priori a derived object. \emph{Caveat:} the
sheaf \eqref{eq:9:general-localization} uses the discrete rings $\mathcal O(U)$,
$\mathcal O^{+}(U)$, so it does not by itself see the topology of a nondiscrete
$\ur R$.

\begin{remark}[It lives over a closed subset]
\label{rmk:9:closed-subset}
In fact $D(\ur R,R^{+})$ lies over a much smaller \emph{closed} subset of the
valuative spectrum,
\begin{equation}
\label{eq:9:closed-subset}
\Spv(\ur R(\ast),R^{+},R^{\circ\circ})\ =\ \bigl\{\,v\in\Spv(\ur R(\ast),R^{+})\
:\ v(f)<1\ \ \forall f\in R^{\circ\circ}\,\bigr\},
\end{equation}
where in addition to $v(f)\le 1$ on $R^{+}$ one imposes $v(f)<1$ strictly for
every topologically nilpotent $f$. For a single $f$ this is a closed condition,
and \eqref{eq:9:closed-subset} is the (closed) intersection over all such $f$.
The claim is that restricting the sheaf of categories to the open complement
gives $0$; so it genuinely lives over \eqref{eq:9:closed-subset}.
\end{remark}

\begin{remark}[Huber's adic spectrum as a retract]
\label{rmk:9:huber-spa}
Huber \cite{Huber_1994} instead works with the \emph{adic spectrum}
\begin{equation}
\label{eq:9:spa}
\Spa(R,R^{+})\ =\ \{\text{continuous valuations on }R,\ \le 1\text{ on }R^{+}\},
\end{equation}
which is a subset of \eqref{eq:9:closed-subset} cut out by a \emph{stronger}
condition: for $f\in R^{\circ\circ}$ and all $\gamma\in\Gamma$ there is
$N\in\mathbb N$ with $v(f^{N})<\gamma$. This subset is generally strictly smaller
than \eqref{eq:9:closed-subset}, the distinction occurring only for higher-rank
valuations --- but in higher dimensions higher-rank valuations are everywhere,
so the difference is large. Now
\begin{itemize}
\item $D(R,R^{+})$ does \emph{not} live over $\Spa(R,R^{+})\subseteq
\Spv(R,R^{+})$ (unlike \eqref{eq:9:closed-subset}, restricting to the complement
does not give $0$);
\item but there is a \emph{retraction}
% board 39--40 @ 01:38:11--01:47:02
\begin{equation}
\label{eq:9:retraction}
\Spv(R,R^{+},R^{\circ\circ})\ \xrightarrow{\ r\ }\ \Spa(R,R^{+}),
\end{equation}
realizing $\Spa(R,R^{+})$ as a \emph{quotient} of the larger spectral space (it
is a priori a subspace), and one \emph{does} get a sheaf of categories on
$\Spa(R,R^{+})$ by pushing forward along $r$. This is the correct way to place a
sheaf of categories on Huber's space.
\end{itemize}
The retraction $r$ is the well-behaved (quasicompact, spectral) map, while the
inclusion $\Spa\hookrightarrow\Spv(\dots,R^{\circ\circ})$ is not spectral; and
$r$ is such that the direct image along it agrees with restriction to the
subspace $\Spa$ (the pertinent direct images vanish, by
\cref{rmk:9:closed-subset}).
\end{remark}

\begin{remark}[More localizations than Huber; the structure sheaf]
\label{rmk:9:more-than-huber}
The valuative-spectrum picture allows strictly more localizations than Huber's.
The rational opens on $\Spa$ correspond to data $f_{1},\dots,f_{n},g$ generating
an \emph{open} ideal; pulling such a rational open back along $r$ recovers the
corresponding rational open upstairs. But a general rational open of
$\Spv(\dots,R^{\circ\circ})$ --- one \emph{not} satisfying the openness
condition --- restricts to something new, not necessarily quasicompact and not a
rational open, which must be written as a union of rational opens. This is like
writing the open unit disk as an increasing union of closed disks; and objects
such as the functions on the closed unit disk (\cref{lec:7}) arise from the
structure sheaf here but not in Huber's setting.

The structure sheaf itself is obtained by taking the globally defined $R$ and
applying the localization functor to land in the category attached to a rational
open; as a symmetric monoidal category the value is generated by its \emph{unit},
which one can analyse. In good cases it corresponds to the Huber pair of the
rational domain, but sometimes only up to a derived enhancement (one must
remember not just the ordinary condensed ring in a full subcategory, but a
derived enhancement of it; the abelian category of modules over the ordinary
thing then suffices). Besides the derived subtlety there is a
quasiseparatedness one: even in degree $0$ the value can differ from Huber's if
the relevant quotient is not by a closed ideal --- though in all practical cases
it does not, and even inverting $g$ can introduce non-quasiseparated behaviour in
general. Iterating the construction shows the global analytic functions of such a
rational domain genuinely carry more data than Huber's (they need not invert a
fixed power of a uniformizer, only converge on each disk), so this theory sees
rings of functions unavailable to the classical formalism.
\end{remark}
